%=======================================================================
%  CHAPTER 2 — Z-TRANSFORM  (4 hrs)
%  Sources: BB Sir decks 2.1, 2.2 (ROC properties, all seven transform
%           properties with proofs); Chapterwise Ch3; NAC Sir CH-3.
%           Cross-checked against Proakis/Manolakis Ch3 and
%           Oppenheim/Schafer Ch3.
%  Numericals live in ch2-num.tex.
%=======================================================================
\section{Z-Transform}

{\color{sub}\footnotesize 4 hrs \gap$\cdot$\gap \mk{61 questions across the 45 papers}
\gap$\cdot$\gap worth about \mk{7 of 80 marks} \gap$\cdot$\gap 3 syllabus sub-topics
\gap$\cdot$\gap the shortest chapter, and \mk{the best marks-per-hour in the subject}.}

\lead{Read this first:}
\begin{itemize}
  \item Only \textbf{two} things are ever asked, and they are asked \emph{together}:
  \begin{itemize}
    \item \textbf{state the properties of the ROC} \tS{14} --- a 1 to 5 mark bookwork
          question, \mk{same answer every year},
    \item \textbf{find the inverse Z-transform} of a given $X(z)$ with a given ROC
          \tS{32}, for 5 or 6 marks.
  \end{itemize}
  \item Those two appear as \textbf{one} question, Q3, in almost every paper. Learn
        \S2.2 and \S2.5 and the chapter is finished.
  \item \mk{The ROC is not decoration.} The same $X(z)$ has a different $x[n]$ for every
        ROC, and papers exploit that: 81 Ba, 82 Ba and 71 Ch each ask for
        \textbf{three} answers from one $X(z)$.
  \item Properties with proofs (convolution \tF{5}, time shift, scaling) are worth 3 to 6
        marks and come from \S2.3.
  \item \mk{Use one method for every inverse}: divide by $z$, partial-fraction
        $X(z)/z$, multiply back so every term reads $z/(z-a)$. \S2.5 explains why.
\end{itemize}

\hr

\T{2.1 Definition, and how it relates to the DTFT}

\Q{Define the Z-transform. How is it related to the DTFT? Explain the scaling
property}{\m{3+3}\m{2}}
{\color{sub}\footnotesize \tP{3} \yr{\textbf{\texttt{74 Bh}}, \textbf{\texttt{76 Bh}},
\texttt{73 Ma}} \gap$\cdot$\gap the definition is the opening line of most Ch 2 answers}

\sbsr{0.52}{%
\lead{Bilateral (two-sided) Z-transform}
\[ X(z)=\sum_{n=-\infty}^{\infty}x[n]\,z^{-n},\qquad z=re^{j\omega} \]
\begin{itemize}
  \item $z$ is a \textbf{complex} variable, so $X(z)$ is defined on a plane, not a line.
  \item The sum exists only where it converges: that region is the \textbf{ROC}, and
        \mk{$X(z)$ without its ROC is only half an answer}.
  \item Splitting the sum shows where the two halves live:
        \[ X(z)=\underbrace{\sum_{n=-\infty}^{-1}x[n]z^{-n}}_{\text{left-sided part}}
                +\underbrace{\sum_{n=0}^{\infty}x[n]z^{-n}}_{\text{right-sided part}} \]
        The right half needs $|z|$ \emph{large}, the left half needs $|z|$ \emph{small};
        an ROC is where both hold at once.
\end{itemize}}{%
\lead{Relation to the DTFT \mk{(quote this)}}
\[ X(e^{j\omega})=X(z)\Big|_{z=e^{j\omega}} \]
\begin{itemize}
  \item Put $r=1$: the Z-transform evaluated \textbf{on the unit circle} is the DTFT.
  \item So the DTFT exists \mk{if and only if the ROC contains the unit circle}. That one
        sentence answers half the theory questions in this chapter.
  \item In reverse: $X(z)$ is the DTFT of the exponentially weighted sequence
        $x[n]r^{-n}$. The $r^{-n}$ is what lets the Z-transform handle growing sequences
        the DTFT cannot.
  \item Going round the unit circle once is $\omega$ advancing by $2\pi$ --- the same
        $2\pi$-periodicity as \S1.4.
\end{itemize}}

\penalty0\vspace{2pt}
\lead{Unilateral (one-sided) Z-transform}
\sbsr{0.42}{%
\[ X^{+}(z)=\sum_{n=0}^{\infty}x[n]\,z^{-n} \]
\begin{itemize}
  \item Ignores $n<0$ entirely, so its ROC is always the outside of a circle and
        \mk{need not be stated}.
  \item \textbf{Importance}: it is the only version that can carry \textbf{initial
        conditions}, so it is what solves a difference equation for a system that is not
        relaxed.
\end{itemize}}{%
\begin{itemize}
  \item \textbf{Time delay property} of the one-sided transform (\texttt{73 Ma} asks this):
        \[ Z^{+}\{x[n-k]\}=z^{-k}X^{+}(z)+\sum_{m=1}^{k}x[-m]\,z^{-(k-m)} \]
  \item \mk{The extra sum is the initial conditions}; the bilateral version has no such
        term, which is exactly why the one-sided transform is the one that solves a
        difference equation with given $y[-1]$, $y[-2]$.
\end{itemize}}

\penalty0\vspace{2pt}
\lead{Standard pairs \mk{(memorise these five)}}
\par
\begin{tabularx}{\textwidth}{@{}L{3.4cm}L{4.0cm}L{2.0cm}X@{}}
\toprule
\hrow \thead{$x[n]$} & \thead{$X(z)$} & \thead{ROC} & \thead{Note} \\
$\delta[n]$ & $1$ & all $z$ & finite duration \\
$u[n]$ & $\dfrac{1}{1-z^{-1}}=\dfrac{z}{z-1}$ & $|z|>1$
 & pole on the unit circle, so \mk{not stable} \\
$a^nu[n]$ & $\dfrac{1}{1-az^{-1}}=\dfrac{z}{z-a}$ & $|z|>|a|$
 & right-sided, causal \\
$-a^nu[-n-1]$ & $\dfrac{1}{1-az^{-1}}=\dfrac{z}{z-a}$ & $|z|<|a|$
 & left-sided, anticausal \\
$n\,a^nu[n]$ & $\dfrac{az^{-1}}{(1-az^{-1})^2}$ & $|z|>|a|$
 & repeated pole; also $\dfrac{az}{(z-a)^2}$ \\
\bottomrule
\end{tabularx}
\par\vspace{2pt}
{\color{sub}\footnotesize \mk{Rows 3 and 4 have the identical $X(z)$ and differ only in the
ROC.} That is the whole reason the ROC must be given, and the reason an inverse cannot be
started until you have read it.}

\hr

\T{2.2 Region of convergence}

\Q{Define the ROC. State and explain the properties of the region of
convergence}{\m{1}\m{2}\m{3}\m{4}\m{5}}
{\color{sub}\footnotesize \tS{14} \yr{\textbf{80 Bh}, \textbf{73 Ch}, \textbf{72 Ch},
\textbf{81 Bh}, \textbf{\texttt{77 Ch}}, \textbf{\texttt{71 Bh}}, \textbf{\texttt{70 Bh}},
79 Ba, 72 Ka, 71 Shr, \textbf{\texttt{73 Bh}}, \textbf{\texttt{78 Ch}}, \texttt{70 Ma},
\texttt{74 Ma}} \gap$\cdot$\gap \mk{pure bookwork, and it is asked every single year}}

\sbsr{0.56}{%
\lead{Definition}
\begin{itemize}
  \item The \textbf{ROC} is the set of values of $z$ for which
        $\sum_n|x[n]||z|^{-n}$ converges, i.e.\ where $X(z)$ actually exists.
  \item It depends only on $|z|=r$, never on the angle, which is why it is always drawn as
        circles.
\end{itemize}
\lead{The eight properties \mk{(this is the answer, in this order)}}
\begin{enumerate}
  \item The ROC is a \textbf{ring or a disc} in the $z$-plane, centred on the origin.
  \item The \textbf{DTFT of $x[n]$ converges absolutely if and only if the ROC includes
        the unit circle.}
  \item \mk{The ROC cannot contain any pole.} $X(z)$ is infinite at a pole, so the sum
        cannot converge there.
  \item If $x[n]$ is \textbf{finite duration}, the ROC is the entire $z$-plane, possibly
        except $z=0$ and/or $z=\infty$.
  \item If $x[n]$ is \textbf{right-sided} ($x[n]=0$ for $n<N_1$), the ROC extends
        \textbf{outward} from the outermost finite pole, possibly including $z=\infty$.
  \item If $x[n]$ is \textbf{left-sided} ($x[n]=0$ for $n>N_2$), the ROC extends
        \textbf{inward} from the innermost non-zero pole, possibly including $z=0$.
  \item If $x[n]$ is \textbf{two-sided}, the ROC is a \textbf{ring} bounded inside and
        outside by a pole, containing none.
  \item The ROC must be a \textbf{connected region}.
\end{enumerate}
{\color{sub}\footnotesize For a 4 or 5 mark answer add one line of \emph{why} to properties
3, 5 and 6, and sketch the three cases. For 1 or 2 marks, the definition plus
properties 1, 3 and 5 is enough.}}{%
\figT{c2_roc_right.png}
\penalty0\vspace{1pt}
{\color{sub}\footnotesize \textbf{Right-sided} $x[n]=a^nu[n]$: the series
$\sum(az^{-1})^n$ converges for $|az^{-1}|<1$, i.e.\ \mk{$|z|>|a|$} --- everything
\emph{outside} the pole circle. \yr{BB Sir, deck 2.1}}
\penalty0\vspace{3pt}
\figT{c2_roc_left.png}
\penalty0\vspace{1pt}
{\color{sub}\footnotesize \textbf{Left-sided} $x[n]=-a^nu[-n-1]$: same
$X(z)=z/(z-a)$, but the ROC is \mk{$|z|<|a|$}, the disc \emph{inside} the pole circle.
\yr{BB Sir, deck 2.1}}}

\sbsr{0.42}{%
\figT{c2_roc_props.png}
\penalty0\vspace{1pt}
{\color{sub}\footnotesize \textbf{Two-sided}: the ROC is the ring between two poles. It
contains the unit circle here, so this signal has a DTFT.
\yr{Chapterwise Ch3}}}{%
\lead{Stability and causality from the ROC \mk{(asked with the properties)}}
\begin{itemize}
  \item \textbf{Causal} $\Leftrightarrow$ $h[n]=0$ for $n<0$ $\Leftrightarrow$ the ROC is
        the \mk{outside of the outermost pole}, and includes $z=\infty$.
  \item \textbf{BIBO stable} $\Leftrightarrow$ $\sum|h[n]|<\infty$ $\Leftrightarrow$ the
        \mk{ROC contains the unit circle}.
  \item \textbf{Causal and stable together} $\Leftrightarrow$ \mk{every pole lies strictly
        inside the unit circle}. This is the single most quoted line in Chapters 2, 3 and 6.
  \item A stable but non-causal system is possible: the ROC is a ring that contains the
        unit circle but does not reach $z=\infty$.
\end{itemize}
\lead{Reading the ROC off a question \mk{(the skill the numericals need)}}
\begin{enumerate}
  \item Find the poles.
  \item Mark their radii on a line.
  \item The given ROC lies in one of the gaps. Any pole with radius \textbf{below} that
        gap is \mk{right-sided}; any pole \textbf{above} it is \mk{left-sided}.
\end{enumerate}
{\color{sub}\footnotesize \mk{A given ROC that contains a pole is impossible} --- two papers
print exactly that, see Ch2-num \S4.}}

\hr

\T{2.3 Properties of the Z-transform}

\Q{State and prove the convolution property; describe causality and stability in terms of
the ROC}{\m{2+4}\m{3+3}\m{1+3+3}}
{\color{sub}\footnotesize \tF{5} \yr{76 Ash, 72 Ka, 70 Asa, \texttt{72 Ma},
\textbf{\texttt{72 Ash}}} \gap$\cdot$\gap \mk{the most-asked property by far}}

\par
\begin{tabularx}{\textwidth}{@{}L{4.0cm}XL{4.6cm}@{}}
\toprule
\hrow \thead{Property} & \thead{Transform} & \thead{ROC} \\
Linearity & $a\,x_1[n]+b\,x_2[n]\;\rightarrow\;aX_1(z)+bX_2(z)$
 & at least $R_{x_1}\cap R_{x_2}$ \\
Time shift & $x[n-k]\;\rightarrow\;z^{-k}X(z)$
 & $R_x$, except possibly $z=0$ or $z=\infty$ \\
Scaling in $z$ (multiply by $a^n$) & $a^n x[n]\;\rightarrow\;X(a^{-1}z)$
 & $|a|r_1<|z|<|a|r_2$ \\
Differentiation & $n\,x[n]\;\rightarrow\;-z\dfrac{dX(z)}{dz}$ & $R_x$ \\
Time reversal & $x[-n]\;\rightarrow\;X(z^{-1})$
 & $\dfrac{1}{r_2}<|z|<\dfrac{1}{r_1}$ \\
Convolution & $x_1[n]*x_2[n]\;\rightarrow\;X_1(z)X_2(z)$
 & at least $R_{x_1}\cap R_{x_2}$ \\
Multiplication & $x_1[n]x_2[n]\;\rightarrow\;\dfrac{1}{2\pi j}\oint_C X_1(v)\,
   X_2\!\left(\dfrac{z}{v}\right)v^{-1}dv$ & at least $r_{1_1}r_{1_2}<|z|<r_{2_1}r_{2_2}$ \\
\bottomrule
\end{tabularx}

\par\vspace{3pt}
\sbs{%
\lead{Convolution, proved \mk{(the 3 to 6 mark answer)}}
\begin{itemize}
  \item Start from the definition of the convolution sum:
        $x[n]=\sum_{k}x_1[k]\,x_2[n-k]$.
  \item Transform both sides:
        \[ X(z)=\sum_{n}\Big[\sum_{k}x_1[k]x_2[n-k]\Big]z^{-n} \]
  \item \textbf{Interchange the order of summation} (the step that earns the mark):
        \[ X(z)=\sum_{k}x_1[k]\sum_{n}x_2[n-k]z^{-n} \]
  \item The inner sum is the \textbf{time-shift property}: $\sum_n x_2[n-k]z^{-n}
        =z^{-k}X_2(z)$.
        \[ X(z)=X_2(z)\sum_{k}x_1[k]z^{-k}=X_1(z)\,X_2(z) \]
  \item \mk{``At least'' the intersection}: a pole of one can be cancelled by a zero of the
        other, which makes the ROC larger than the intersection. Say this --- it is the
        part most answers miss.
\end{itemize}}{%
\lead{Time shift, proved \m{3} \yr{\textbf{69 Ch}}}
\begin{itemize}
  \item Let $y[n]=x[n-k]$, so
        $Y(z)=\sum_n x[n-k]z^{-n}$.
  \item Substitute $m=n-k$, so $n=m+k$:
        \[ Y(z)=\sum_m x[m]z^{-(m+k)}=z^{-k}\sum_m x[m]z^{-m}=z^{-k}X(z) \]
  \item \mk{Why the ROC can change}: $z^{-k}$ adds poles at $z=0$ (for $k>0$) or at
        $z=\infty$ (for $k<0$), so those two points may drop out.
\end{itemize}
\lead{Scaling, proved \m{3} \yr{\textbf{\texttt{74 Bh}}}}
\begin{itemize}
  \item[] \[ Z\{a^nx[n]\}=\sum_n a^nx[n]z^{-n}=\sum_n x[n](a^{-1}z)^{-n}=X(a^{-1}z) \]
  \item The ROC $r_1<|z|<r_2$ becomes $r_1<|a^{-1}z|<r_2$, i.e.\
        $|a|r_1<|z|<|a|r_2$: \mk{every pole is scaled by $a$}.
  \item Example: $u[n]\rightarrow\frac{1}{1-z^{-1}}$, so
        $a^nu[n]\rightarrow\frac{1}{1-az^{-1}}$, pole moved from 1 to $a$.
\end{itemize}
\lead{Differentiation, used in every $na^n$ question}
\begin{itemize}
  \item Differentiate $X(z)=\sum x[n]z^{-n}$ term by term and multiply by $-z$:
        \[ Z\{n\,x[n]\}=-z\frac{dX(z)}{dz} \]
  \item Applied to $a^nu[n]$ it gives
        $Z\{na^nu[n]\}=\dfrac{az^{-1}}{(1-az^{-1})^2}$, $|z|>|a|$ ---
        \mk{the answer \texttt{78 Ch} and \texttt{70 Ma} both ask for}.
\end{itemize}}

\hr

\T{2.4 Finding $X(z)$ and locating the ROC}

\Q{Find the Z-transform of the given signal and locate its ROC}{\m{2+4}\m{3+6}\m{4+6}}
{\color{sub}\footnotesize \tF{6} \yr{\textbf{70 Ch}, 70 Asa, 75 Ash, \textbf{75 Ch},
\textbf{\texttt{75 Bh}}, \texttt{73 Ma}} \gap$\cdot$\gap worked in \textbf{Ch2-num \S6}}

\sbs{%
\lead{Method}
\begin{enumerate}
  \item Split $x[n]$ into its terms.
  \item Transform each with a standard pair, \mk{writing its own ROC beside it}.
  \begin{itemize}
    \item $a^nu[n]\rightarrow\frac{1}{1-az^{-1}}$, ROC $|z|>|a|$ (outside).
    \item $-b^nu[-n-1]\rightarrow\frac{1}{1-bz^{-1}}$, ROC $|z|<|b|$ (inside).
    \item \mk{Watch the sign}: a $+b^nu[-n-1]$ term transforms to
          $-\frac{1}{1-bz^{-1}}$.
  \end{itemize}
  \item The ROC of the sum is the \textbf{intersection} of the individual ROCs.
  \item If the intersection is empty, \mk{the Z-transform does not exist} --- and that is
        the full answer. 70 Ch prints exactly such a signal (Ch2-num \S6).
\end{enumerate}}{%
\lead{Why the intersection can be empty}
\begin{itemize}
  \item A right-sided term needs $|z|>R_1$; a left-sided term needs $|z|<R_2$.
  \item The ring $R_1<|z|<R_2$ exists \mk{only if $R_1<R_2$}: the right-sided pole must be
        the \emph{inner} one.
  \item $R_1=R_2$ or $R_1>R_2$ leaves nothing, so $X(z)$ does not exist even though each
        term separately transforms perfectly well.
\end{itemize}}

\penalty0\vspace{2pt}
\lead{Two standard transforms the papers ask for by name}
\sbs{%
\begin{itemize}
  \item $x[n]=na^nu[n]$ (\textbf{\texttt{78 Ch}}, \texttt{70 Ma}) --- use
        differentiation:
        \[ X(z)=\frac{az^{-1}}{(1-az^{-1})^2},\qquad |z|>|a| \]
  \item Derivation in one line: $-z\dfrac{d}{dz}\dfrac{z}{z-a}
        =-z\cdot\dfrac{(z-a)-z}{(z-a)^2}=\dfrac{az}{(z-a)^2}$.
\end{itemize}}{%
\begin{itemize}
  \item $x[n]=\cos\omega n$ for $n\ge0$ (\texttt{74 Ma}) --- write the cosine as two
        exponentials and add two geometric series:
        \[ X(z)=\frac{1-z^{-1}\cos\omega}{1-2z^{-1}\cos\omega+z^{-2}},\qquad |z|>1 \]
        \mk{Its two poles are $e^{\pm j\omega}$, on the unit circle}, which is why the
        ROC excludes it and why a pure sinusoid is not BIBO stable.
\end{itemize}}

\penalty0\vspace{1pt}

\hr

\T{2.5 Inverse Z-transform: the three methods}

\Q{Determine the inverse Z-transform of the given $X(z)$ for the stated
ROC}{\m{5}\m{6}\m{1+5}\m{2+6}}
{\color{sub}\footnotesize \tS{32} \gap$\cdot$\gap \mk{the single biggest question in this
chapter}, 5 or 6 marks in nearly every paper \gap$\cdot$\gap every instance is worked in
\textbf{Ch2-num}}

\sbs{%
\lead{Which method, and when}
\begin{enumerate}
  \item \textbf{Partial fractions} --- the default. Use unless the question names another
        method. Gives a closed form for every $n$.
  \item \textbf{Long division (power series)} --- use when the question says so
        (\textbf{\texttt{77 Ch}}, \textbf{\texttt{76 Bh}}, \textbf{\texttt{73 Bh}},
        \texttt{72 Ma}, \texttt{70 Ma}). Gives numbers, \mk{not a closed form}, so give
        4 or 5 terms and the recurrence.
  \item \textbf{By inspection} --- when $X(z)$ is already a finite polynomial, read the
        coefficients straight off as impulses (74 Ash, \textbf{\texttt{78 Ch}}).
  \item The contour-integral / residue method is never asked in these 45 papers.
\end{enumerate}}{%
\lead{\mk{Why divide by $z$ first}}
\begin{itemize}
  \item The standard pairs are $\dfrac{z}{z-a}\leftrightarrow a^nu[n]$, with a $z$
        \emph{on top}.
  \item Expanding $X(z)$ directly gives terms $\dfrac{A}{z-a}$ with no $z$ on top, and
        each then needs an extra time shift. Expanding $X(z)/z$ puts the $z$ back for
        free.
  \item \mk{Use it every single time}, including the repeated-pole cases. One idiom, no
        decisions.
\end{itemize}}

\penalty0\vspace{2pt}
\sbs{%
\lead{The partial-fraction method, in one fixed order}
\begin{enumerate}
  \item \textbf{Is $X(z)$ proper?} If the numerator's degree is $\ge$ the denominator's,
        \mk{long-divide first} and keep the polynomial part aside --- it becomes impulses.
  \item \textbf{Divide by $z$} and expand $X(z)/z$ in partial fractions.
  \item Multiply back by $z$, so every term reads $\dfrac{A\,z}{z-a}$.
  \item \textbf{Assign each pole a side} from the ROC (\S2.2).
  \item Write the answer term by term with the pairs opposite.
\end{enumerate}}{%
\lead{The pair table, both sides}
\par
\begin{tabular}{@{}L{2.4cm}L{2.9cm}L{2.6cm}@{}}
\toprule
\hrow \thead{Term} & \thead{ROC $|z|>|a|$} & \thead{ROC $|z|<|a|$} \\
$\dfrac{z}{z-a}$ & $a^n u[n]$ & $-a^n u[-n-1]$ \\
$\dfrac{z}{(z-a)^2}$ & $n\,a^{n-1}u[n]$ & $-n\,a^{n-1}u[-n-1]$ \\
constant $C$ & $C\delta[n]$ & $C\delta[n]$ \\
$z^k$ & \multicolumn{2}{l}{$\delta[n+k]$, an \emph{advance}} \\
\bottomrule
\end{tabular}
\par\vspace{2pt}
{\color{sub}\footnotesize \mk{The right-hand column is just the left with a minus sign and
$u[-n-1]$ instead of $u[n]$.} That is the only thing the ROC changes.}}

\lead{Long division: which way to divide \mk{(where the marks are lost)}}
\sbs{%
\begin{itemize}
  \item \textbf{Right-sided} (ROC outside, causal): arrange both polynomials in
        \mk{ascending powers of $z^{-1}$} and divide. The quotient
        $x[0]+x[1]z^{-1}+x[2]z^{-2}+\ldots$ is read off directly.
  \item \textbf{Left-sided} (ROC inside, anticausal): arrange in \mk{ascending powers of
        $z$} instead. The quotient gives $x[-1]z,\,x[-2]z^2,\ldots$
  \item \mk{Dividing the wrong way round gives the other sequence}, which is a valid
        series for a different ROC --- and scores zero.
\end{itemize}}{%
\begin{itemize}
  \item Faster than dividing: read the \textbf{recurrence} straight off the denominator.
        For $X(z)=\dfrac{1}{1-a_1z^{-1}-a_2z^{-2}}$,
        \[ x[n]=a_1x[n-1]+a_2x[n-2],\quad x[0]=1 \]
        and then just iterate. This is the same difference equation the filter implements
        (Ch 4).
  \item \mk{Always give the recurrence as well as the terms}: it shows the pattern
        continues and it is what a 6-mark answer wants.
\end{itemize}}

\hr

\T{2.6 Last-minute recall for this chapter}

\sbs{%
\lead{Must memorise}
\begin{itemize}
  \item $X(z)=\sum_n x[n]z^{-n}$; the DTFT is $X(z)$ on the unit circle.
  \item $a^nu[n]\leftrightarrow\frac{z}{z-a}$, $|z|>|a|$; \;
        $-a^nu[-n-1]\leftrightarrow\frac{z}{z-a}$, $|z|<|a|$. \mk{Same $X(z)$, different
        ROC.}
  \item $na^nu[n]\leftrightarrow\frac{az^{-1}}{(1-az^{-1})^2}$;
        $\frac{z}{(z-a)^2}\leftrightarrow na^{n-1}u[n]$.
  \item The 8 ROC properties (\S2.2) --- \mk{written out, in order}.
  \item Causal $\Rightarrow$ ROC outside the outermost pole. Stable $\Rightarrow$ ROC
        contains the unit circle. \mk{Both $\Rightarrow$ all poles inside the unit circle.}
  \item Convolution $\rightarrow$ product; shift $\rightarrow$ $z^{-k}$; multiply by $a^n$
        $\rightarrow$ $X(a^{-1}z)$; multiply by $n$ $\rightarrow$ $-z\,dX/dz$.
  \item Inverse: proper? $\rightarrow$ divide by $z$ $\rightarrow$ partial fractions
        $\rightarrow$ assign sides from the ROC.
\end{itemize}}{%
\lead{Most repeated, in order}
\begin{enumerate}
  \item Inverse Z-transform with a given ROC \tS{32}
  \item Properties of the ROC \tS{14}
  \item Find $X(z)$ and locate its ROC \tF{6}
  \item Convolution property, stated and proved \tF{5}
  \item Z-transform of $na^nu[n]$ \tP{3}
  \item One-sided Z-transform and its time-delay property \tP{2}
  \item Time-shift property, derived \tP{1}
  \item Scaling property with an example \tP{1}
\end{enumerate}
{\color{sub}\footnotesize \mk{Q3 of the paper is the ROC properties plus an inverse
Z-transform in 30 of the 45 papers.} If you prepare one question in this subject, prepare
that one.}}
